% 第4章 放大电路的频率响应
\chapter{放大电路的频率响应}

\section{频率响应概述}
\subsection{研究放大电路频率响应的必要性}
实际输入信号由不同频率的正弦波组成，放大电路对不同频率信号的放大能力不同（增益、相移不同），若频率响应不佳会导致信号失真，因此需分析频率响应。

\subsection{频率响应的基本概念}
\subsubsection{\texorpdfstring{高通电路}{高通电路}}
高通电路允许高频信号通过，抑制低频信号，存在**下限截止频率** $f_L$（下限截止频率处增益下降3dB）。
典型高通电路：耦合电容、发射极旁路电容构成的电路。

\subsubsection{\texorpdfstring{低通电路}{低通电路}}
低通电路允许低频信号通过，抑制高频信号，存在**上限截止频率** $f_H$（上限截止频率处增益下降3dB）。
典型低通电路：晶体管极间电容、分布电容构成的电路。

\subsubsection{\texorpdfstring{波特图}{波特图}}
将放大电路的频率响应（增益-频率、相移-频率）用对数坐标绘制的图形，分为：
\begin{itemize}
    \item 幅频特性波特图：纵坐标为增益的分贝数 $20\lg|A_u|$，横坐标为频率的对数 $\lg f$；
    \item 相频特性波特图：纵坐标为相移 $\varphi$（单位：°），横坐标为 $\lg f$。
\end{itemize}
通频带 $f_{BW} = f_H - f_L$，表示放大电路能有效放大的信号频率范围。

\section{\texorpdfstring{晶体管的高频等效模型}{晶体管的高频等效模型}}
\subsection{\texorpdfstring{晶体管的混合$\pi$模型}{晶体管的混合pi模型}}
为分析晶体管高频特性，将晶体管等效为混合$\pi$结构，核心参数：
\begin{itemize}
    \item $r_{bb'}$：基区体电阻；
    \item $r_{b'e}$：发射结等效电阻，$r_{b'e} = (1+\beta_0) \frac{V_T}{I_{EQ}}$；
    \item $C_{b'e}$：发射结结电容；
    \item $C_{b'c}$：集电结结电容（米勒电容）；
    \item $g_m$：跨导，$g_m = \frac{I_{CQ}}{V_T}$（$V_T \approx 26\ \text{mV}$，室温）。
\end{itemize}

\subsection{\texorpdfstring{晶体管电流放大倍数的频率响应}{晶体管电流放大倍数的频率响应}}
晶体管共射电流放大倍数 $\beta$ 随频率升高而下降，满足：
\[ \beta = \frac{\beta_0}{1 + j\frac{f}{f_\beta}} \]
其中：
\begin{itemize}
    \item $\beta_0$：低频段电流放大倍数；
    \item $f_\beta$：$\beta$ 的截止频率（$\beta$ 下降到 $\beta_0/\sqrt{2}$ 时的频率）；
    \item 特征频率 $f_T$：$|\beta|=1$ 时的频率，$f_T \approx \beta_0 f_\beta$（高频重要参数）。
\end{itemize}

\section{\texorpdfstring{场效应管的高频等效模型}{场效应管的高频等效模型}}
场效应管高频等效模型（类似晶体管混合$\pi$模型）核心参数：
\begin{itemize}
    \item $r_{gs}$：栅源极间电阻（阻值极大，可近似开路）；
    \item $C_{gs}$、$C_{gd}$：栅源、栅漏极间电容；
    \item $g_m$：跨导，$g_m = \frac{\partial I_D}{\partial U_{GS}} \big|_{U_{DS}=\text{常数}}$；
    \item $r_{ds}$：漏源极间输出电阻。
\end{itemize}
场效应管极间电容更小，高频特性优于晶体管。

\section{\texorpdfstring{单管放大电路的频率响应}{单管放大电路的频率响应}}
\subsection{\texorpdfstring{单管共射放大电路的频率响应}{单管共射放大电路的频率响应}}
单管共射放大电路的频率响应分为三个频段：
\begin{enumerate}
    \item 中频段：电容（耦合电容、旁路电容）视为短路，极间电容视为开路，增益恒定、无相移；
    \item 低频段：耦合/旁路电容的容抗增大，导致增益下降、产生超前相移，下限截止频率 $f_L$ 由耦合/旁路电容决定；
    \item 高频段：极间电容的容抗减小，导致增益下降、产生滞后相移，上限截止频率 $f_H$ 由晶体管混合$\pi$模型的电容决定。
\end{enumerate}
单管共射放大电路的全频段电压放大倍数：
\[ A_{us} = \frac{A_{usm}}{\left(1 - j\frac{f_L}{f}\right)\left(1 + j\frac{f}{f_H}\right)} \]
其中 $A_{usm}$ 为中频段源电压放大倍数。

\section{\texorpdfstring{多级放大电路的频率响应}{多级放大电路的频率响应}}
\subsection{\texorpdfstring{多级放大电路频率响应的特点}{多级放大电路频率响应的特点}}
1. 总电压放大倍数：$A_u = \prod_{k=1}^n A_{uk}$（各级放大倍数乘积）；
2. 总幅频特性：$20\lg|A_u| = \sum_{k=1}^n 20\lg|A_{uk}|$（各级分贝数相加）；
3. 通频带：多级放大电路的通频带比单级窄，级数越多，通频带越窄；
4. 截止频率：
   - 总下限截止频率 $f_L$：大于各级下限截止频率，$f_L \approx 1.1\sqrt{f_{L1}^2 + f_{L2}^2 + \dots + f_{Ln}^2}$；
   - 总上限截止频率 $f_H$：小于各级上限截止频率，$\frac{1}{f_H} \approx 1.1\sqrt{\frac{1}{f_{H1}^2} + \frac{1}{f_{H2}^2} + \dots + \frac{1}{f_{Hn}^2}}$。

\section{例题参考}
5